\subsection{Tight convergence of bounded measures}

Let \(\mathbf{T}\) be a completely regular space; the bilinear form

\[
(f, \mu) \mapsto \int f (t) d \mu (t)
\]

on \(\mathcal{C}^{b}(\mathrm{T})\times\mathcal{M}^{b}(\mathrm{T})\) puts these two spaces in a separating duality. For, it is clear that the duality is separating in \(\mathcal{C}^{b}(\mathrm{T})\) from the fact that the measures \(\varepsilon_{x}\;(x\in\mathrm{T})\) belong to \(\mathcal{M}^{b}(\mathrm{T})\) ; it is separating in \(\mathcal{M}^{b}(\mathrm{T})\) by Prop. 2 of No.~1.

DEFINITION 1. --- The weak topology on \(\mathcal{M}^{b}(\mathrm{T})\) associated with the preceding duality between \(\mathcal{C}^{b}(\mathrm{T})\) and \(\mathcal{M}^{b}(\mathrm{T})\) is called the topology of tight convergence (or the tight topology) on \(\mathcal{M}^{b}(\mathrm{T})\) .

The tight topology is Hausdorff, by the remarks preceding the definition. We shall often employ the adverb `tightly' to mean `in the sense of the tight topology'. Absent mention to the contrary, \(\mathcal{M}^b (\mathrm{T})\) will be equipped with the tight topology throughout the rest of this section.

Every element of \(\mathcal{C}^{b}(\mathrm{T})\) is a linear combination of elements of \(\mathcal{C}_{+}^{b}(\mathrm{T})\) . For a filter F on \(\mathcal{M}^{b}(\mathrm{T})\) to converge tightly to a bounded measure \(\lambda\) , it is necessary and sufficient that

\begin{enumerate}
\def\labelenumi{(\arabic{enumi})}
\setcounter{enumi}{3}
\tightlist
\item
\end{enumerate}

\(\lim_{\mu} \mu(f) = \lambda(f)\) with respect to \(\mathfrak{F}\) for every \(f \in \mathcal{C}_+^b(\mathbf{T})\) .

Remarks. --- 1) If T is locally compact, the tight topology is finer than the topology induced on \(\mathcal{M}^{b}(\mathrm{T})\) by the vague topology, and these two topologies coincide only when T is compact. For, if T is not compact, the mapping \(t \mapsto \varepsilon_{t}\) converges vaguely to 0 with respect to the filter of complements of relatively compact subsets of T, but does not converge tightly to 0, because the function 1 belongs to \(\mathcal{C}^{b}(\mathrm{T})\) (for the relations between vague convergence and tight convergence, see Prop. 9).

\begin{enumerate}
\def\labelenumi{\arabic{enumi})}
\setcounter{enumi}{1}
\item
  It follows at once from Prop. 4 that \(\mathcal{M}_+^b (\mathbf{T})\) is closed in \(\mathcal{M}^b (\mathbf{T})\) .
\item
  If \(\mathbf{T}\) is completely regular, the mapping \(t \mapsto \varepsilon_t\) of \(\mathbf{T}\) into \(\mathcal{M}^b(\mathbf{T})\) is a homeomorphism (GT, IX, §1, No.~5).
\end{enumerate}

PROPOSITION 6. --- Let T be a completely regular space.

\begin{enumerate}
\def\labelenumi{\alph{enumi})}
\item
  Let \(f\) be a lower semi-continuous numerical function \(\geqslant 0\) defined on \(\mathbf{T}\) ; then the function \(\mu \mapsto |\mu|^{\bullet}(f)\) is lower semi-continuous on \(\mathcal{M}^b(\mathbf{T})\) .
\item
  Let \(f\) be an upper semi-continuous bounded function defined on \(\mathbf{T}\) ; then the function \(\mu \mapsto \mu(f)\) is upper semi-continuous on \(\mathcal{M}_+^b(\mathbf{T})\) .
\end{enumerate}

For, one sees by Prop. 1 b) of No.~1 that \(\mu \mapsto |\mu|^{\bullet}(f)\) is the upper envelope of a family of functions of the form \(\mu \mapsto |\mu(g)|\) with \(g \in \mathcal{C}^b(T)\) , hence continuous for the tight topology. This establishes \(a\) ). To prove \(b\) ), it suffices to choose a constant upper bound C for \(f\) , and to write \(\mu(f) = \mu(C) - \mu(C - f)\) ; the function \(\mu \mapsto \mu(C)\) is continuous, and the function \(\mu \mapsto \mu(C - f)\) is lower semi-continuous on \(\mathcal{M}_+^b(T)\) by the foregoing.

PROPOSITION 7. --- Let T be a completely regular space. Let \(\mu\) be a bounded positive measure on T, and let f be a bounded positive function on T, such that the set of points of T where f is not continuous is locally \(\mu\) -negligible. Then the mapping \(\lambda \mapsto \lambda^{\bullet}(f)\) of \(\mathcal{M}_{+}^{b}(\mathrm{T})\) into R is continuous at the point \(\mu\) .

For every \(t \in \mathbb{T}\) , set \(f'(t) = \liminf_{s \to t} f(s)\) , \(f''(t) = \limsup_{s \to t} f(s)\) . Obviously \(f' \leqslant f \leqslant f''\) , with equality at every point of \(\mathbb{T}\) where \(f\) is continuous (hence \(\mu\) -almost everywhere). On the other hand, \(f'\) is lower semi-continuous, \(f''\) is upper semi-continuous and bounded (GT, IV, §6, No.~2, Prop. 4). We therefore have the following relations by Prop. 6,

\[
\begin{array}{c} \mu^ {\bullet} (f ^ {\prime}) \leqslant \liminf _ {\lambda \to \mu} \lambda^ {\bullet} (f ^ {\prime}) \leqslant \liminf _ {\lambda \to \mu} \lambda^ {\bullet} (f) \leqslant \mu^ {\bullet} (f) \leqslant \limsup _ {\lambda \to \mu} \lambda^ {\bullet} (f) \\ \leqslant \limsup _ {\lambda \to \mu} \lambda^ {\bullet} (f ^ {\prime \prime}) \leqslant \mu^ {\bullet} (f ^ {\prime \prime}). \end{array}
\]

One concludes by observing that \(\mu^{\bullet}(f') = \mu^{\bullet}(f'')\) , because \(f'\) and \(f''\) are equal locally \(\mu\) -almost everywhere.

PROPOSITION 8. --- Let X be a completely regular space, T a subspace of X, and i the canonical injection of T into X. Denote by W the set of bounded positive measures on X that are concentrated on T, equipped with the topology induced by \(\mathcal{M}^{b}(X)\) . Then the mapping \(\mu \mapsto i(\mu)\) of \(\mathcal{M}_{+}^{b}(T)\) into \(\mathcal{M}^{b}(X)\) is a homeomorphism of \(\mathcal{M}_{+}^{b}(T)\) onto W.

We denote again by \(i\) the mapping \(\mu \mapsto i(\mu)\) of \(\mathcal{M}_{+}^{b}(\mathrm{T})\) into \(\mathcal{M}_{+}^{b}(\mathrm{X})\) ; \(i\) is injective (§2, No.~4, Prop. 8) and maps \(\mathcal{M}_{+}^{b}(\mathrm{T})\) into W (§2, No.~3, Prop. 7). If \(\lambda \in \mathrm{W}\) , then \(\lambda = i(\lambda_{\mathrm{T}})\) (§2, No.~3, Prop. 7 b)). Consequently, \(i\) is a bijection of \(\mathcal{M}_{+}^{b}(\mathrm{T})\) onto W, and the inverse bijection of \(i\) is the mapping \(r : \lambda \mapsto \lambda_{T}\) on W. On the other hand, i is continuous: for, if \(f \in \mathcal{C}^{b}(X)\) , then \(\langle i(\mu), f \rangle = \langle \mu, f \circ i \rangle\) , and \(f \circ i\) belongs to \(\mathcal{C}^{b}(T)\) . Thus, everything comes down to showing that, for every measure \(\mu \in W\) and every function \(f \in \mathcal{C}_{+}^{b}(T)\) , one has

\[
\lim _ {\lambda \to \mu , \lambda \in W} \lambda_ {T} (f) = \mu_ {T} (f),
\]

or again

\[
\lim _ {\lambda \to \mu ,   \lambda \in \mathrm{W}} \lambda (f ^ {0}) = \mu (f ^ {0}).
\]

Let \(f^{\infty}\) be the function on X that coincides with f on T and with \(+\infty\) on X - T, and let \(f'\) and \(f''\) be, respectively, the upper semi-continuous regularization of \(f^{0}\) and the lower semi-continuous regularization of \(f^{\infty}\) (GT, IV, §6, No.~2). The relations

\[
f ^ {\prime} (x) = \limsup _ {y \to x} f ^ {0} (y), \quad f ^ {\prime \prime} (x) = \liminf _ {y \to x} f ^ {\infty} (y)
\]

immediately imply that \(f'\) and \(f''\) both coincide with f and \(f^{0}\) on T. Prop. 6 then yields

\[
\mu^ {\bullet} (f ^ {\prime}) \geqslant \lim _ {\lambda \to \mu ,   \lambda \in W} \sup \lambda^ {\bullet} (f ^ {\prime})  , \quad \mu^ {\bullet} (f ^ {\prime \prime}) \leqslant \lim _ {\lambda \to \mu ,   \lambda \in W} \inf \lambda^ {\bullet} (f ^ {\prime \prime})  .
\]

But one can replace \(f'\) and \(f''\) by \(f^{0}\) in these two formulas, since the measures \(\lambda\) and \(\mu\) are carried by T; we have thus obtained the desired relation.

The statement of Prop. 8 is only valid for positive measures: the mapping \(\mu \mapsto i(\mu)\) of \(\mathcal{M}^b (\mathrm{T})\) into \(\mathcal{M}^b (\mathrm{X})\) is injective and continuous, but is not in general a homeomorphism of \(\mathcal{M}^b (\mathrm{T})\) onto its image. For example, take \(\mathbf{X} = \mathbf{R}\) , \(\mathrm{T} = \mathbf{R} - \{0\}\) ; the measures \(\lambda_t = \varepsilon_t - \varepsilon_{-t}\) ( \(t > 0\) ) converge tightly to 0 in X as \(t\) tends to 0, but do not converge tightly to 0 in T (the characteristic function of {]}0, +∞{[} belongs to \(\mathcal{C}^b (\mathrm{T})\) ) (cf.~however the Cor. of Th. 1 of No.~5).

PROPOSITION 9. --- Let T be a locally compact space, and let F be a filter on \(\mathcal{M}_{+}^{b}(T)\) that converges vaguely to a bounded measure \(\mu\) . For F to converge tightly to \(\mu\) , it is necessary and sufficient that \(\lim_{\lambda}\lambda(1)=\mu(1)\) with respect to F.

The condition is obviously necessary. To show that it is sufficient, let us denote by X the Alexandroff compactification of T (GT, I, §9, No.~8) and by i the canonical injection of T into X. By Prop. 8, everything comes down to showing that \(\lambda \mapsto i(\lambda)\) converges tightly to \(i(\mu)\) in \(\mathcal{M}^{b}(X)\) with respect to F. Since \(\mu(1) < +\infty\) , there exists a set \(A \in F\) such that the total masses of the measures in A are bounded by a number M; it therefore suffices to verify that

\[
\lim _ {\lambda , \mathfrak {F}} \int_ {\mathrm{X}} g d (i (\lambda)) = \int_ {\mathrm{X}} g d (i (\mu))\tag{5}
\]

for functions \(g \in \mathcal{C}^b(\mathrm{X})\) forming a total set in \(\mathcal{C}^b(\mathrm{X})\) . Now, this equality is satisfied when \(g\) has compact support in \(\mathrm{T}\) , because of the vague convergence of \(\mathfrak{F}\) to \(\mu\) , and also when \(g\) is a constant function on \(\mathrm{X}\) , from the fact that \(\lim_{\lambda, \mathfrak{F}} \lambda(1) = \mu(1)\) . Since the functions of the preceding two types form a total set in \(\mathcal{C}^b(\mathrm{X})\) (Ch. III, §1, No.~2, Prop. 3), this completes the proof.

